\chapter{Downtime Management}
\label{ch:downtime}
\index{downtime}\index{GM!downtime}

Downtime is the narrative space between adventures -- days, weeks, or even months when characters rest, train, build assets, settle debts, and pursue personal projects. \textbf{As GM, your role is to:}  
\begin{itemize}
  \item Set duration and DVs based on fiction, not XP budgets.
  \item Hear player declarations and call for appropriate rolls.
  \item Apply the Outcome Matrix to determine progress/complications.
  \item Spend Story Beats from player rolls to introduce rumors, rival actions, or unexpected events.
  \item Manage campaign timers (Mandate, Crisis, faction) based on player choices.
  \item Track Asset/Follower upkeep and enforce neglect consequences.
  \item Run the Echo Character system for off-screen PCs.
  \item Adjudicate Rush Rules for accelerated development.
\end{itemize}
Unlike many games, Fate's Edge treats downtime as an active phase where \textbf{your adjudication shapes the fiction}--not a mere pause. This chapter gives you the tools to run downtime efficiently and dramatically.

\section{GM Procedures for Running Downtime}
\label{sec:gm-downtime-procedures}
\index{downtime!procedures}\index{GM!downtime}
Each downtime phase follows this loop. \textbf{Follow these steps as GM:}  

\begin{enumerate}
  \item \textbf{Set the Duration.} 
        Determine how much time passes based on fiction (usually 1--4 weeks for standard play; longer for epic campaigns). 
        \textit{Example:} ``After the siege, two weeks pass while the army regroups--what do you do with this time?''
  \item \textbf{Hear Player Declarations.} 
        Each player states one or two major goals (training, research, carousing, etc.). 
        \textit{Note:} Players declare intent--you set DVs and call for rolls.
  \item \textbf{Call for Rolls \& Resolve.} 
        For each declaration: 
        \begin{itemize}
          \item Set a Difficulty Value (DV) using the Common Downtime Rolls table (see \ref{sec:downtime-ref}) as a guide; increase by +1/+2 for rushed/dangerous projects.
          \item Have the player roll Attribute + Skill d10s; count Successes (6+) and Story Beats (1s).
          \item Apply the Outcome Matrix (Clean Success / Success with SB / Partial / Miss).
          \item \textbf{Note all Story Beats generated} for later spending.
        \end{itemize}
  \item \textbf{Spend Story Beats.} 
        Use any Story Beats generated during downtime to: 
        \begin{itemize}
          \item Introduce a minor rumor (1 SB): ``A street urchin whispers you're being watched by the Black Ledger.''
          \item Signal faction awareness (2 SB): ``The City Guard captain mentions your name at roll call.''
          \item Advance a complication timer (3+ SB): ``The \emph{Harbormaster's Grudge [4]} timer ticks +1--2.''
        \end{itemize}
        \textit{GM Tip:} Spend SB to create \textit{immediate hooks} for the next session--never waste them.
  \item \textbf{Manage Upkeep \& Timers.} 
        \begin{itemize}
          \item \textbf{Asset/Follower Upkeep}: Remind players to pay XP or devote downtime scenes. If neglected, degrade condition (Maintained $\rightarrow$ Neglected $\rightarrow$ Compromised) and apply consequences (see \ref{sec:followers-assets-gm}).
          \item \textbf{Campaign Timers}: Advance Mandate/Crisis/faction timers based on player choices (see \ref{sec:downtime-timers}).
          \item \textbf{Obligation Clearing}: When players clear Obligation via acts of service, describe the narrative cost (e.g., ``You spend the day repairing the temple roof--your hands blister, but the priests nod in approval'').
        \end{itemize}
  \item \textbf{Handle Off-Screen Characters.} 
        Process Downtime Logs using the Echo Character system (see \ref{sec:echo-system}).
        \textit{GM Tip:} Treat off-screen characters as active NPCs--their logs are your source for complications and opportunities.
  \item \textbf{Log the Results.} 
        Use the downtime log template to record key outcomes for session-to-session continuity.
\end{enumerate}

\section{Common Downtime Rolls Reference (GM)}
\label{sec:downtime-ref}
\index{downtime!reference}
Use this table to set DVs for player declarations. Increase DV by +1 (risky) or +2 (dangerous) as fiction demands.

\renewcommand{\arraystretch}{1.2}
\begin{longtable}{p{4cm}p{2.5cm}p{3cm}p{5cm}}
\toprule
\textbf{Player Goal} & \textbf{Skill Roll} & \textbf{Base DV} & \textbf{Effect on Success} \\
\midrule
\endfirsthead
\multicolumn{4}{c}{{\bfseries Continued from previous page}} \\
\toprule
\textbf{Player Goal} & \textbf{Skill Roll} & \textbf{Base DV} & \textbf{Effect on Success} \\
\midrule
\endhead
\bottomrule
\endfoot
Recover from Harm & Medicine + Wits & 3 & Remove 1 Harm (stacks with rest) \\
Train a Skill & relevant Attribute + Skill & 4 & Advance 1 segment on a Training Timer [4]. \\
Learn a New Talent & appropriate Attribute + Lore & 5 & Advance 1 segment on a Talent Acquisition Timer [6]. \\
Clear Obligation & Spirit + Lore (act of service) & 3 & Reduce Obligation by 1 segment (you describe the act). \\
Carouse for Information & Presence + Streetwise & 3 & Gain a rumor or a minor String. \\
Craft an Item & Wits + Craft & DV by item (3--6) & Create a mundane or minor magical item. \\
Research a Topic & Wits + Investigation & 4 & Gain one piece of true information (you may be cryptic). \\
Manage an Asset & Presence + Command & 3 & Reduce or prevent condition degradation. \\
Establish a Contact & Presence + Sway & 4 & Gain a new Minor String. \\
Reduce a Debt Timer & Presence + Sway or service roll & 4 & Clear 1--2 segments (your discretion). \\
Train a Follower & Spirit + Command & 4 & Improve Follower Cap by 1 (requires Training Timer). \\
\end{longtable}

\section{Rush Rules: GM Procedures for Risk Timers}
\label{sec:rush-rules}
\index{Rush rule}\index{Risk Timer}\index{GM!rush}
Players can attempt to rush downtime activities. \textbf{As GM, follow this procedure:}  

\begin{enumerate}
  \item \textbf{Player declares rushing** (e.g., ``I want to learn Acrobatics in 3 days instead of a week''). 
  \item \textbf{You create a Risk Timer} of 4--6 segments, named after the potential consequence (e.g., \textit{Pulled Muscle [6]}, \textit{Backlash [4]}, \textit{Patron Ire [4]}). 
        \textit{Narrate the risk:} ``Pushing this hard could tear a tendon--I'm starting a \textit{Pulled Muscle [6]} timer.''
  \item \textbf{For each downtime action related to the rush, the player rolls as normal.} 
        \begin{itemize}
          \item On a \textbf{Clean Success}: The timer does not advance. 
          \item On a \textbf{Partial} or \textbf{Miss}: The timer advances by 1 segment (or more, at your discretion for extreme failure). 
        \end{itemize}
  \item \textbf{When the Risk Timer fills}, choose one outcome based on fiction: 
        \begin{itemize}
          \item \textbf{Project fails entirely**, but the player gains 2 Boons to use on a later attempt. 
          \item \textbf{Project succeeds, but the character suffers a permanent Condition** (e.g., \textit{Twisted Ankle}, \textit{Migraines}) that lasts until treated. 
          \item \textbf{Character gains the new ability, but immediately marks 1 Harm (overexertion) and gains a Debt Timer [4]** to a terrestrial patron who ``helped'' them cheat time. 
        \end{itemize}
  \item \textbf{If the timer never fills**, the rush succeeds as normal with no extra cost. 
\end{enumerate}

\begin{tcolorbox}[colback=red!5!white, colframe=red!75!black, title=Example: Rushing a Skill (GM Narration),breakable]
Michael wants to raise his Melee from 3 to 4. Normal training takes 4 days. He rushes to do it in 2 days. You create a \textit{Pulled Muscle [6]} timer: ``Push too hard, and you'll tear something.'' Michael rolls three downtime actions (one per day). He gets a Partial and a Miss -- the timer fills. You rule: ``You gain the skill, but during the final drill you tear a tendon. Mark Harm 1 (Level 1) and you have the \textit{Limp} Condition until you receive magical healing or rest for a full week.''
\end{tcolorbox}

\section{Using Downtime to Advance Campaign Timers}
\label{sec:downtime-timers}
\index{timers!downtime}\index{GM!timers}
Downtime lets players actively shape campaign threats and favor. \textbf{As GM, adjudicate these declarations:}  

\subsection{Mandate and Crisis}
\label{subsec:mandate-crisis}
Players can declare they are working to increase Mandate (public favor) or decrease Crisis (threat level). 
\begin{itemize}
  \item \textbf{Increase Mandate (public favor):} 
        Player declares charity, problem-solving, or alliance-building. 
        \textbf{You set DV 4} for Presence + Sway. 
        \begin{itemize}
          \item On success: Increase Mandate by 1 (to a max of 6). 
          \item On a Miss: Increase Crisis by 1 (enemies notice their growing influence). 
        \end{itemize}
  \item \textbf{Decrease Crisis (reduce threat):} 
        Player declares sabotage, truce-brokering, or conspiracy exposure. 
        \textbf{You set DV 5} for an appropriate skill (e.g., Stealth, Subterfuge, Lore). 
        \begin{itemize}
          \item On success: Reduce Crisis by 1. 
          \item On a Miss: Mandate also decreases by 1 (collateral damage--you describe how). 
        \end{itemize}
\end{itemize}

\subsection{Faction and Debt Timers}
\label{subsec:faction-debt-timers}
Players can influence visible timers through downtime actions. 
\begin{itemize}
  \item \textbf{Faction Timers}: 
        Player declares action to help or harm a faction (e.g., ``I bribe officials in the Dyers' Guild''). 
        \textbf{You set DV 4} for Presence + Sway. 
        \begin{itemize}
          \item On success: Tick the faction timer \textbf{back} by 1 (less threat). 
          \item On a Miss: Tick the timer \textbf{forward} by 1 (more threat). 
          \item \textit{Example:} Success on bribing officials: ``Your coins find the right palms--the Guild Influence timer retreats one segment.''
        \end{itemize}
  \item \textbf{Debt Timers from Terrestrial Patrons}: 
        Player accepts a Debt Timer during downtime (see \ref{sec:obligation-corruption-gm}). 
        \textbf{You resolve it via}: 
        \begin{itemize}
          \item A service scene (roll to clear 1--2 segments), 
          \item Or spending XP (1 XP = 1 segment cleared). 
        \end{itemize}
\end{itemize}

\section{Handling Off-Screen Characters: The Echo System}
\label{sec:echo-system}
\index{Echo character}\index{off-screen}\index{GM!echo}
When a player has multiple characters or steps a character off-screen, \textbf{you use this system to keep them active in the fiction}:  

\subsection{Core Principles}
\begin{itemize}
  \item Off-screen characters are \textbf{not frozen**. They continue to act, gather information, and face consequences. 
  \item The player provides a brief \textbf{Downtime Log} for each off-screen character (you provide the template). 
  \item You spend 2--3 minutes per log, describing outcomes and possibly ticking a personal \textbf{Crisis Timer [6]} for that character. 
\end{itemize}

\subsection{Switch Mechanics}
\label{subsec:switch-mechanics}
When a player wants to bring an off-screen character back into play (or send one off-screen): 
\begin{enumerate}
  \item \textbf{Declare} -- Player states: ``I'm switching to my other character, the scholar, for the next session.'' 
  \item \textbf{Status} -- \textbf{You assign** the current character one of three statuses based on fiction: 
        \begin{itemize}
          \item \textbf{Supporting} -- Works toward active PC goals (gathers intel, manages assets). Minimal risk. 
          \item \textbf{Advancing} -- Pursues their own agenda (research, training). May uncover opportunities or trigger a complication. 
          \item \textbf{In Peril} -- Their absence has made them a target. \textbf{You tick a personal Crisis Timer [6]} each session until they return. 
        \end{itemize}
  \item \textbf{Re-entry} -- When returning, the character gains 2 Boons for a dramatic entrance. 
        The player rolls Wits + Spirit (DV 3) to re-establish presence. 
        \textbf{You adjudicate the roll}: 
        \begin{itemize}
          \item On a Miss: A complication follows them back (you describe it using 1--2 SB). 
        \end{itemize}
\end{enumerate}

\subsection{Downtime Log Template (Provide to Players)}
\label{sec:downtime-log}
\index{downtime log}\index{GM!log}
Give players this simple log to fill out for each off-screen (or on-screen) character between sessions:  

\begin{tcolorbox}[colback=gray!5!white, colframe=black, title=Downtime Log -- \_\_\_\_\_\_\_\_\_ (Character Name),breakable]
\begin{tabular}{p{0.9\textwidth}}
\textbf{Resource Action:} (How they maintain connections, pay upkeep, or keep assets running) \\
\rule{0pt}{4ex}\\
\textbf{Project Action:} (One goal they worked toward -- training, research, crafting) \\
\rule{0pt}{4ex}\\
\textbf{Event Response:} (How they react to a rumor, threat, or off-screen event you describe) \\
\rule{0pt}{4ex}\\
\textbf{GM Response:} (You fill this in after the session) \\
\rule{0pt}{4ex}
\end{tabular}
\end{tcolorbox}

\subsection{Off-Screen Crisis Timer }
Each off-screen character tracks a personal \textbf{Crisis Timer [6]}. \textbf{You tick this timer when}:  
\begin{itemize}
  \item The player rolls a Miss on a downtime action. 
  \item The character's status is \textit{In Peril}. 
  \item You spend a Story Beat to introduce trouble for that character. 
\end{itemize}
When the timer fills, \textbf{you introduce a major off-screen consequence}: a rival strikes, an asset is destroyed, or they are captured. This creates immediate hooks for the active party. 

\section{Additional Downtime Tools and Advice (GM)}
\label{sec:downtime-advice}
\index{downtime!advice}\index{GM!advice}
Use these tools to make downtime a tension-filled narrative phase:  

\subsection{Asset \& Follower Upkeep}
Remind players to pay upkeep (Efficient or Intensive) for each asset and follower. Use the templates from Chapter 10. If they forget, \textbf{degrade the asset to Neglected or a follower to Wary}. You can also use the Rush rules to restore a degraded asset in half the normal downtime. 

\subsection{Resolving Complications}
Downtime is an excellent moment to clear lingering Complications. Players can spend a downtime action to resolve a complication--roll appropriate skill (DV 3--5). \textbf{As GM, adjudicate:} 
\begin{itemize}
  \item On success: The complication is cleared (describe the narrative resolution). 
  \item On a Miss: It worsens--\textbf{you advance a related timer by 1 segment}. 
\end{itemize}

\subsection{Using Story Beats from Downtime Rolls}
Any 1s rolled during downtime generate Story Beats. \textbf{Do not waste them. Spend them as follows}: 
\begin{itemize}
  \item \textbf{1 SB:} A minor rumor begins about the character (``A beggar claims they saw you arguing with the watch captain''). 
  \item \textbf{2 SB:} A rival or faction becomes aware of their activities (``The City Guard captain mentions your name at roll call''). 
  \item \textbf{3 SB:} A complication arises--a theft, a broken item, a social faux pas (``Your safehouse door is jammed--spend 1 Boon to force it open''). 
  \item \textbf{4+ SB:} A major twist--a Patron demands an audience, an enemy strikes an asset (``The Harbormaster confronts you at dawn: 'Your games end today.'''). 
\end{itemize}
These complications will ignite the next session's opening scene. 

\subsection{Downtime as Collaborative Worldbuilding}
Encourage players to ask questions during downtime (``I try to find a fence who deals in cursed coins. Is there one in this city?''). \textbf{Use their rolls to establish new NPCs and locations}: 
\begin{itemize}
  \item \textbf{Clean Success}: The fence exists and is reliable (``Old Jem's stall in the market--he'll take any cursed silver for fair price''). 
  \item \textbf{Partial}: The fence exists but has a price (Debt Timer)--you describe the cost (``Old Jem nods, but his eyes flick to your purse: 'Five silver upfront, or no deal.'''). 
  \item \textbf{Miss}: They exist but are working for an enemy (``Old Jem's stall is abandoned--in its place, a Black Ledger enforcer eyes you coldly''). 
\end{itemize}

\subsection{Sample Downtime Sequence (GM Script)}
Use this template to frame downtime at your table:  
\begin{quote}
\textquotedbl [Time period] passes. You each have [number] downtime actions. I need: a roll for skill training, a roll for asset upkeep, and a roll for whatever personal project you declared. You can also use a downtime action to interact with the campaign timers--tell me how you want to affect the Mandate or the Cult's Progress timer. Describe your goal, I'll set a DV, and we'll roll. Any 1s will give me Story Beats to complicate your return to action. Ready?\textquotedbl
\end{quote}

\section{Downtime Quick Reference (GM Cheat Sheet)}
\label{sec:downtime-cheat-sheet}
\begin{tcolorbox}[colback=blue!5!white, colframe=blue!75!black, title=Downtime GM Cheat Sheet,breakable]
\begin{itemize}
  \item \textbf{Set duration} based on fiction (1--4 weeks typical; adjust for scene weight). 
  \item \textbf{Hear player declarations} (1--2 goals per player). 
  \item \textbf{Set DVs} using reference table; increase for risk/rush. 
  \item \textbf{Call for rolls}; players count Successes/SBs. 
  \item \textbf{Apply Outcome Matrix}; note SBs generated. 
  \item \textbf{Spend SB} to introduce rumors/rival actions/complications (1 SB = minor rumor; 2 SB = faction awareness; 3 SB = complication; 4+ SB = major twist). 
  \item \textbf{Advance/retreat campaign timers} (Mandate, Crisis, faction) based on player choices + fiction. 
  \item \textbf{Process off-screen characters}: Collect Downtime Logs, apply Echo system mechanics. 
  \item \textbf{Enforce Rush Rule}: Create Risk Timer [4--6] for rushed actions; process Partial/Miss as timer ticks. 
  \item \textbf{Track Asset/Follower upkeep**: Players report; you enforce degradation if neglected. 
  \item \textbf{Resolve complications**: Player rolls vs DV; success = cleared, miss = worsens (you advance timer). 
\end{itemize}
\end{tcolorbox}
\end{document}
